\documentclass[output=paper]{langscibook}
\ChapterDOI{10.5281/zenodo.22880260}
\title{\textit{A kiu vaa teo atovo} – The work on the sago palm thatch}
\author{Enoch Horai Magum}
\abstract{}
\begin{document}
\maketitle\label{chap:17}
\begin{figure}[H]
    \captionsetup[subfigure]{oneside,margin={0cm,0cm}}
    \subfloat[]{\includegraphics[trim=0 112 0 30,clip,width=.4\linewidth]{figures/17_1.jpg}\label{fig:17_1}}\hfil
    \subfloat[]{\includegraphics[trim=0 0 0 0,clip,width=.6\linewidth]{figures/17_2.jpg}\label{fig:17_2}}
    \caption{(a) \textit{O Atovo} -- The Sago palm, (b) \textit{Paka atovo} -- Sago palm leaves}
\end{figure}
\begin{xltabular}{\textwidth}{lXX}
    (\ref{ex:text12-1})&\textit{O atovo eve o meho taba to pakupaku inu ni raara en Teapu.}&The sago palm thatch is a thing that we use to build houses here in the Teop region.\\
    (\ref{ex:text12-2})&\textit{Eve to vavaakopo riori toon na inu tea varavihi kahi bona huan, a roava ae a iravata. A kiu vaa teo atovo he na paku vonaen:}&It is put on top of the house to protect it against the rain, the sun and the storm. The work on the sago palm thatch is done as follows:\\
    (\ref{ex:text12-3})&\textit{Beara  suguna teo hum to tei nae a atovo, eara toro viana vamataa o vuaba vai o manoto vamataa, bara to vapepeana vamataa nana mee.}&When we have arrived at the place where the sago palms are, we must carefully choose one that is fully mature and also has many branches.\\
    (\ref{ex:text12-4})&\textit{Beara viana vakavara, eara repaa tagava e ge eara repaa peha e.}&When we have chosen (it), we fell it or we climb it.\\
    (\ref{ex:text12-5})&\textit{Tea peha, eara na peha rara, eara repaa mosi koa maa a peana gina beo atovo sun vakisi pahi nana.}&As for the climbing, we climb it, and then we only cut the branches on the top, while the sago palm is still standing.\\
    (\ref{ex:text12-6})&\textit{Tea tagava, eara na tagava rara o atovo, ore obete nana putaa beara mosi rara a peana.}&As for the felling, we fell the sago palm so that it lies on the ground, while we are cutting off the branches.\\
    (\ref{ex:text12-7})&\textit{Tea mosi peana he eara mee toro viana batara, are goe tea gono a peana vai to saka tara vamataa haana vai, a peana a tasitasiri ge to tara vakobusu nana.}&As for the cutting of branches, we must also be selective, so that we won’t get the branches that do not look good, branches that are torn or that look young and  moist.\\
    (\ref{ex:text12-8})&\textit{Beara mosi vakavara a pea atovo, eara nao vai tea kiu vaa tea kosi; eara kou kosi vai a paka atovo kahi bona pea nae.}&When we have cut the sago palm branches, we now go on with the (next) work of cutting off the leaves; we now cut the sago-palm leaves off their branch.\\
    (\ref{ex:text12-9})&\textit{Be kavara o hum bona, eara nao vai tea kiu vaa tea varovaromo nia paka atovo.}&When this part is finished, we go on with the work of stacking up the sago palm leaves.\\
    (\ref{ex:text12-10})&\textit{Tea varovaromo he eara na varomo ni rara a ta taonim ge a ta peha totoka paka atovo.}&When stacking them up, we stack up five or six sago palm leaves.\\
    (\ref{ex:text12-11})&\textit{Bea kiu vaa tea varovaromo kavara, eara kou paku vai a peha dee i. Eara repaa kisi e, kisi rara bona kukan a kisi, eara repaa dee voosu.}&When the work of stacking up is finished, we make one load. We tie it up by tying up three bundles, and then we carry it home.\\
    (\ref{ex:text12-12})&\textit{A peha aba na antee nana tea kosi bono dorana o peana.}&A single person can cut off the leaves of four branches.\\
    (\ref{ex:text12-13})&\textit{Ei beara suguna vaan, eara paku pete vai a meha kiu, a kiu vaa tea siri.}&Now, when we have arrived in the village, we do the next work, the work of tearing.\\
\end{xltabular}
\begin{figure}[H]
    \captionsetup[subfigure]{oneside,margin={0cm,0cm}}
    \subfloat[]{\includegraphics[trim=0 0 0 165,clip,width=.38\linewidth]{figures/17_3.jpg}\label{fig:17_3}}
    \subfloat[]{\includegraphics[trim=0 0 0 0,clip,width=.62\linewidth]{figures/17_4.jpg}\label{fig:17_4}}
    \caption{(a) \textit{Siri} -- Tearing, (b) \textit{Kaakaku, pisi} -- Breaking, folding}
\end{figure}
\begin{xltabular}{\textwidth}{lXX}
    (\ref{ex:text12-14})&\textit{Tea siri, eara na siri kahi batara a paka no atovo bona sikiri nae, ore tamee vamataa beara huha batara.}&As for the tearing, we tear the midrib off the sago palm leaf, so that it (the leaf) is flexible when we are sewing (the thatch sheets).\\
    (\ref{ex:text12-15})&\textit{Beara siri vakavara, eara repaa nao pete tea kiu vaa tea varovaromo vahabana nie, eara repaa kaakaku e.}&When we have finished the tearing, we go on with the work of stacking them up again, and then we fold them.\\
    (\ref{ex:text12-16})&\textit{Evehee eara toro kaku murina batara to peho ipinoo.}&But we must fold them according to one measure.\\
    (\ref{ex:text12-17})&\textit{Bea kaakaku kavara, eara repaa vaobete vatataunun e, eara repaa kisi vahabana nie.}&When the folding is finished, we put them together, one on top of the other, then we tie them up again.\\
\end{xltabular}
\begin{figure}[H]
    \captionsetup[subfigure]{oneside,margin={0cm,0cm}}
    \subfloat[]{\includegraphics[trim=0 60 0 0,clip,width=.5\linewidth]{figures/17_5.jpg}\label{fig:17_5}}\hfil
    \subfloat[]{\includegraphics[trim=0 0 0 0,clip,width=.5\linewidth]{figures/17_6.jpg}\label{fig:17_6}}
    \caption{(a) \textit{Vaobete vatataunun} -- Putting together, (b) \textit{Amaa kisi atovo} -- Bundles of sago palm leaves}
\end{figure}
\begin{xltabular}{\textwidth}{lXX}
    (\ref{ex:text12-18})&\textit{Eara repaa vaobete vamataa e teo hum to raka vamataa nana.}&We carefully lay them in a place that is completely dry.\\
    (\ref{ex:text12-19})&\textit{Pasi obete nom teebona, beara ihuana rara o bon vaa tea huha.}&They will lie there, while we are waiting for the sewing day.\\
    (\ref{ex:text12-20})&\textit{Eara na antee rara tea pahin huha bata ei bea tataraha matatopo nana ge ei bea naba kiu bero rori.}&We can now directly sew (the thatch), when the \textit{kururu} bamboo strip (on which the leaves are sewn)\footnote{See \textit{tataraha}\textsuperscript{2} in the \textsc{mted} \url{https://dictionaria.clld.org/units/teop-tataraha_2}} is prepared, or when the group of workers is big.\\
\end{xltabular}
\begin{figure}[H]
    \captionsetup[subfigure]{oneside,margin={0cm,0cm}}
    \subfloat[]{\includegraphics[trim=0 0 0 8,clip,width=.41\linewidth]{figures/17_7.jpg}\label{fig:17_7}}
    \subfloat[]{\includegraphics[trim=0 0 0 80,clip,width=.27\linewidth]{figures/17_8.jpg}\label{fig:17_8}}
    \subfloat[]{\includegraphics[trim=0 0 0 15,clip,width=.32\linewidth]{figures/17_9.jpg}\label{fig:17_9}}
    \caption{(a) \textit{Kururu} bamboo, (b) and (c) \textit{Huha atovo} -- Sewing a sago palm thatch sheet}
\end{figure}
\begin{xltabular}{\textwidth}{lXX}
    (\ref{ex:text12-21})&\textit{Evehee a bero amaa vuri toro obete roho.}&But many days the (bundles) must lie (there) before (we can start).\\
    (\ref{ex:text12-22})&\textit{Tea kiu vaa tea huha he na paku vonaen:}&The work of sewing, however, is done as follows:\\
    (\ref{ex:text12-23})&\textit{Eara na vaobete rara a tataraha toon no kasuana, eara repaa parata o paka atovo, ei kou na kaku pisi momohu ara.}&We put the bamboo strip on the ground, then we open a sago-palm leaf, because we had folded it before.\\
    (\ref{ex:text12-24})&\textit{Erau eara repaa vaobete maa a tataraha komano paka atovo.}&Then we lay the \textit{kururu} bamboo strip inside the (folded) sago palm leaf.\\
    (\ref{ex:text12-25})&\textit{Eara repaa vaakopo taono vahae bona.}&Then we cover it up again.\\
    (\ref{ex:text12-26})&\textit{Eara repaa huu nio koopu toon na ta atovo.}&Then we stick the \textit{koopu} bamboo (needle) from above into the upper part of the (folded) sago-palm thatch sheet.\\
    (\ref{ex:text12-27})&\textit{Eara repaa tahi bono sikiri vai to siri kahi ara bono atovo teo vuri siri atovo.}&Then we insert the midrib that we have torn off the sago palm leaf when doing the sago palm leaf tearing.\\
    (\ref{ex:text12-28})&\textit{Eara na tahi hopo rara, eara repaa tahi tavusu maa,}&We push it into (the hole of the folded leaf), and then we push it out.\\
    (\ref{ex:text12-29})&\textit{Eara repaa vaobete vahabana nao o meho paka, eara repaa huha bata naora.}&Then we lay another leaf on top and go on sewing.\\
    (\ref{ex:text12-30})&\textit{Eara pasi paku bata koa vo naora teebona, ore tavusu komana to haha nae a ta tataraha.}&We will do it continuously just like this, until the strip is used up. \\
    (\ref{ex:text12-31})&\textit{Teo nin na vaapeha koana.}&With a \textit{nin} reed, it is just the same.\footnote{See \textit{nin}\textsuperscript{2} in the \textsc{mted} \url{https://dictionaria.clld.org/units/teop-nin_2}}\\
    (\ref{ex:text12-32})&\textit{Evehee ahiki kou to koopu, eara na tahi bata koa ni naora o nin,}&If there is no \textit{koopu} bamboo, we insert a \textit{nin} reed,\\
    (\ref{ex:text12-33})&\textit{ei kou o tamee, bara eara na vavakakaniva rara o hum tahi.}&because it is flexible, and we carve a sharp pointed end that we push into the leaf. \\
    (\ref{ex:text12-34})&\textit{Be kavara a kiu vaa tea huha, eara repaa paku pete a kiu vaa tea vaave.}&When the sewing is finished, we do the next work of putting up the thatch.\\
\end{xltabular}
\begin{figure}[H]
    \captionsetup[subfigure]{oneside,margin={0cm,0cm}}
    \subfloat[]{\includegraphics[trim=50 13 50 0,clip,width=.33\linewidth]{figures/17_10.jpg}\label{fig:17_10}}\hfil
    \subfloat[]{\includegraphics[trim=5 10 20 0,clip,width=.33\linewidth]{figures/17_11.jpg}\label{fig:17_11}}\hfil
    \subfloat[]{\includegraphics[trim=0 0 0 40,clip,width=.34\linewidth]{figures/17_12.jpg}\label{fig:17_12}}
    \caption{(a) \textit{A ta atovo} -- Sago palm thatch sheet, (b) \textit{Vaave} -- Thatching, (c) \textit{A atovo} -- The thatch}
\end{figure}
\begin{xltabular}{\textwidth}{lXX}
    (\ref{ex:text12-35})&\textit{Tea vaave, eara  na paku roho rara a poka, o hum to sun bata raara beara vaave batara.}&Before putting up the thatch, we make a scaffold, a place on which we stand, when we are putting up the thatch.\\
    (\ref{ex:text12-36})&\textit{Me tea vaave, eara na antee rara tea vaave hiava ge vaave gunaha.}&And when thatching, we can thatch upwards (starting from the bottom) or thatch downwards (starting from the ridgepole).\\
    (\ref{ex:text12-37})&\textit{A buaku a aba pasi paku bata nom bona kiu vaa tea hee hiava ni bona atovo, bea meha papana vaave batari.}&Two people would do the work of reaching up the thatch sheets, while others would fix them. \\
    (\ref{ex:text12-38})&\textit{Eara na kisi ni rara o kara sareai, o nin ge o obai, teo varuvaru.}&We tie (a thatch sheet) onto the rafter with a string made of the \textit{sareai} vine,\footnote{We could not identify its English or scientific name of \textit{sareai}; but see Chapter \ref{chap:11}, \figref{fig:11_3}.} the \textit{nin} reed, or the calamus cane.\footnote{See Chapter \ref{chap:11}:\ref{ex:text7-44}ff.}\\
    (\ref{ex:text12-39})&\textit{Bea ta atovo obete maa tea varuvaru, eara repaa kisi taviri ie bona varuvaru.}&When the thatch sheet lies on the rafters, we tie it (the string) around the rafters.\\
    (\ref{ex:text12-40})&\textit{Be kavara, eara repaa dee hiava vahabana a ta meha, ore suguna komana to antee nae a varuvaru.}&When it is finished, we reach up another (thatch sheet), until the end of the rafters is reached.\\
    (\ref{ex:text12-41})&\textit{Tea vaave hiava, eara na vaobete taono bata naora a ta atovo bona ta meha, ore no suguna komana teo vatetekoio.}&When the thatching moves upwards, we put one sheet of sago palm thatch on top of the other, until the ridgepole is reached.\\
    (\ref{ex:text12-42})&\textit{Bea vaave gunaha, eara na taneo maara teo vatetekoio, ore gunaha maa beara hogo bata ni naora a ta atovo paana ta meha.}&When the thatching moves downwards, we start at the ridgepole, and then it goes down, when we shove one sago palm thatch sheet underneath the other.\\
    (\ref{ex:text12-43})&\textit{Bea peha papa inu kavara, eara repaa taneo a meha papa inu.}&When one side of the house is finished, we start with the other side.\\
    (\ref{ex:text12-44})&\textit{Be kavara, ei kavara vai, eara repaa vaobete maa o kovuaka tea mahan vai to paa vatatana nom a vaave, ei teo vatetekoio.}&When (this other side) is finished -- this (thatching) is finished now, and we put the top thatch sheet over the open space where the (two sides of the) thatch meet, this is on the upper ridge pole.\\
\end{xltabular}
\end{document}